% bismi ALLAH ARRAHMAN ARRAHIM




























































%%%%%%%%%%%%%%%%%%%%%%%%%%%%%%%%%%%%%%%%%%%%%%%%%%%%%%%%%%%%%%%%%%%%%%%%%%
\section{Architectures d'apprentissage guidées par l'oracle variationnel}
%%%%%%%%%%%%%%%%%%%%%%%%%%%%%%%%%%%%%%%%%%%%%%%%%%%%%%%%%%%%%%%%%%%%%%%%%%

Le routage variationnel développé dans les sections précédentes fournit une
loi de décision exacte

\[
c^\star(w)
\in
\arg\min_{c\in\mathcal C}
\widetilde V_c(w).
\]

Cette loi n'est pas apprise.

Elle est entièrement déterminée par le problème variationnel local.

Elle constitue donc un \emph{oracle de routage}.

L'idée développée dans cette section est que cet oracle peut également servir
de mécanisme d'apprentissage.

Autrement dit, le même objet mathématique joue simultanément deux rôles :

\begin{enumerate}
\item décider du contexte optimal ;
\item produire automatiquement les données nécessaires à l'apprentissage
d'un routeur rapide.
\end{enumerate}

Aucune base de données annotée n'est alors nécessaire.

Les exemples d'apprentissage sont générés pendant le fonctionnement même
du système.

%%%%%%%%%%%%%%%%%%%%%%%%%%%%%%%%%%%%%%%%%%%%%%%%%%%%%%%%%%%%%%%%%%%%%%%%%%

\subsection{Oracle variationnel}

\begin{definition}[Oracle variationnel]

On appelle oracle variationnel l'application

\[
\mathcal O :
\mathbb R^d
\rightarrow
\mathcal C
\]

définie par

\[
\boxed{
\mathcal O(w)
=
c^\star(w)
=
\arg\min_{c\in\mathcal C}
\widetilde V_c(w).
}
\]

\end{definition}

L'oracle est entièrement déterminé par les problèmes variationnels locaux.

Il ne possède aucun paramètre entraînable.

Ses décisions sont donc entièrement explicables.

%%%%%%%%%%%%%%%%%%%%%%%%%%%%%%%%%%%%%%%%%%%%%%%%%%%%%%%%%%%%%%%%%%%%%%%%%%

\subsection{Réseau de routage}

On considère un réseau

\[
G_\theta :
\mathbb R^d
\longrightarrow
\Delta_K,
\]

où

\[
G_\theta(w)
=
(p_1(w),\ldots,p_K(w))
\]

désigne une distribution de probabilité sur les contextes.

La décision du réseau est

\[
\widehat c_\theta(w)
=
\arg\max_c p_c(w).
\]

%%%%%%%%%%%%%%%%%%%%%%%%%%%%%%%%%%%%%%%%%%%%%%%%%%%%%%%%%%%%%%%%%%%%%%%%%%

\subsection{Apprentissage auto-supervisé}

Chaque appel à l'oracle produit automatiquement

\[
(w,\mathcal O(w)).
\]

Cette paire constitue un exemple d'apprentissage.

Le réseau est mis à jour en minimisant

\[
L(\theta)
=
-
\log
p_{\mathcal O(w)}(w).
\]

Ainsi,

aucune annotation extérieure n'est requise.

L'ensemble d'apprentissage est construit en temps réel.

%%%%%%%%%%%%%%%%%%%%%%%%%%%%%%%%%%%%%%%%%%%%%%%%%%%%%%%%%%%%%%%%%%%%%%%%%%

\begin{proposition}

L'apprentissage précédent est entièrement auto-supervisé.

Les étiquettes ne proviennent pas d'une intervention humaine mais du problème
variationnel lui-même.

\end{proposition}

%%%%%%%%%%%%%%%%%%%%%%%%%%%%%%%%%%%%%%%%%%%%%%%%%%%%%%%%%%%%%%%%%%%%%%%%%%

\subsection{Architecture générale}

Toute architecture hybride peut être décrite par une politique

\[
\Pi
:
\mathbb R^d
\times
\Theta
\longrightarrow
\{0,1\},
\]

où

\[
\Pi(w,\theta)=
\begin{cases}
1,&\text{appel de l'oracle},\\
0,&\text{utilisation directe du réseau}.
\end{cases}
\]

Le routage final est

\[
R(w)
=
\begin{cases}
\mathcal O(w),
&
\Pi(w,\theta)=1,
\\
\widehat c_\theta(w),
&
\Pi(w,\theta)=0.
\end{cases}
\]

Toutes les variantes ci-dessous correspondent simplement à différents choix
de la politique
\(
\Pi.
\)

%%%%%%%%%%%%%%%%%%%%%%%%%%%%%%%%%%%%%%%%%%%%%%%%%%%%%%%%%%%%%%%%%%%%%%%%%%

\subsection{Variante I : oracle permanent}

La politique est

\[
\Pi(w,\theta)\equiv1.
\]

Le réseau apprend en permanence tandis que les décisions restent celles de
l'oracle.

Cette variante fournit la référence théorique.

%%%%%%%%%%%%%%%%%%%%%%%%%%%%%%%%%%%%%%%%%%%%%%%%%%%%%%%%%%%%%%%%%%%%%%%%%%

\subsection{Variante II : alternance oracle--réseau}

On fixe un seuil

\[
\tau\in(0,1).
\]

Si

\[
\max_c p_c(w)\ge\tau,
\]

le réseau décide.

Sinon,

l'oracle est appelé.

Autrement dit,

\[
\Pi(w,\theta)
=
\mathbf1
\left(
\max_c p_c(w)<\tau
\right).
\]

%%%%%%%%%%%%%%%%%%%%%%%%%%%%%%%%%%%%%%%%%%%%%%%%%%%%%%%%%%%%%%%%%%%%%%%%%%

\subsection{Variante III : oracle fondé sur l'entropie}

On définit

\[
H(w)
=
-
\sum_c
p_c(w)
\log p_c(w).
\]

Pour un seuil

\[
\gamma>0,
\]

on choisit

\[
\Pi(w,\theta)
=
\mathbf1
(H(w)>\gamma).
\]

L'oracle n'est donc interrogé que lorsque
l'incertitude est élevée.

%%%%%%%%%%%%%%%%%%%%%%%%%%%%%%%%%%%%%%%%%%%%%%%%%%%%%%%%%%%%%%%%%%%%%%%%%%

\subsection{Variante IV : oracle intermittent}

On fixe

\[
q\in(0,1].
\]

À chaque observation,

l'oracle est interrogé avec probabilité

\[
q.
\]

Ainsi,

\[
\mathbb P
(\Pi=1)
=
q.
\]

Cette stratégie réduit le coût de calcul tout en maintenant une correction
régulière du réseau.

%%%%%%%%%%%%%%%%%%%%%%%%%%%%%%%%%%%%%%%%%%%%%%%%%%%%%%%%%%%%%%%%%%%%%%%%%%

\subsection{Variante V : distillation variationnelle}

Au lieu de transmettre uniquement

\[
c^\star,
\]

l'oracle transmet la distribution

\[
p^\star_\beta(c|w)
=
\frac
{
e^{-\beta\widetilde V_c(w)}
}
{
\sum_j
e^{-\beta\widetilde V_j(w)}
}.
\]

Le réseau minimise

\[
L(\theta)
=
D_{KL}
\left(
p^\star_\beta
\|
G_\theta
\right).
\]

Le réseau apprend alors la géométrie complète des valeurs réduites.

%%%%%%%%%%%%%%%%%%%%%%%%%%%%%%%%%%%%%%%%%%%%%%%%%%%%%%%%%%%%%%%%%%%%%%%%%%

\subsection{Variante VI : pré-routage appris}

Le réseau fournit

\[
S(w)
\subset
\mathcal C,
\]

avec

\[
|S(w)|=m\ll K.
\]

L'oracle est ensuite restreint à

\[
S(w),
\]

et calcule

\[
\arg\min_{c\in S(w)}
\widetilde V_c(w).
\]

Le coût du routage devient proportionnel à

\[
m
\]

au lieu de

\[
K.
\]

%%%%%%%%%%%%%%%%%%%%%%%%%%%%%%%%%%%%%%%%%%%%%%%%%%%%%%%%%%%%%%%%%%%%%%%%%%

\subsection{Variante VII : apprentissage permanent}

Après chaque appel de l'oracle,

les paramètres sont actualisés selon

\[
\theta_{n+1}
=
\theta_n
-
\eta
\nabla_\theta
L
\left(
G_{\theta_n}(w_n),
\mathcal O(w_n)
\right).
\]

Le réseau converge progressivement vers la loi de décision induite par
l'oracle.

%%%%%%%%%%%%%%%%%%%%%%%%%%%%%%%%%%%%%%%%%%%%%%%%%%%%%%%%%%%%%%%%%%%%%%%%%%

\begin{remark}

Toutes les variantes précédentes reposent sur le même principe fondamental.

Le routage variationnel constitue une vérité mathématique produite par la
minimisation des valeurs réduites.

Le réseau n'apprend jamais directement le problème de routage.

Il apprend uniquement à approximer une loi de décision déjà définie par
l'analyse variationnelle.

Le rôle de l'apprentissage est donc exclusivement celui d'une accélération
algorithmique.

La correction théorique demeure assurée par l'oracle variationnel.

\end{remark}

%%%%%%%%%%%%%%%%%%%%%%%%%%%%%%%%%%%%%%%%%%%%%%%%%%%%%%%%%%%%%%%%%%%%%%%%%%

\begin{theorem}[Principe général des architectures guidées par PCH]

Soit

\[
\mathcal O(w)
=
\arg\min_{c}
\widetilde V_c(w)
\]

l'oracle variationnel.

Toute architecture définie par une politique
d'interrogation

\[
\Pi
\]

vérifie les propriétés suivantes.

\begin{enumerate}

\item
Le routage exact est retrouvé lorsque
\(
\Pi\equiv1.
\)

\item
Le réseau construit automatiquement ses données
d'entraînement.

\item
Aucune annotation extérieure n'est requise.

\item
Le coût asymptotique diminue au fur et à mesure que
le réseau remplace progressivement l'oracle.

\item
Lorsque l'oracle demeure accessible,
la correction variationnelle du système reste garantie.

\end{enumerate}

\end{theorem}
